\begin{table}[htbp]
\centering
\small
\caption{Problema A (linear) versus Problema D (meia-lua) (openai/gpt-4o-mini), por regime.}
\label{tab:bloco1_A_vs_D_por_regime_gpt4mini}
\begin{tabular}{lrrrrrr}
\toprule
 & \multicolumn{2}{c}{Fidelidade (vs LLM)} & \multicolumn{2}{c}{Acur.\ da métrica vs real} & \multicolumn{2}{c}{Acur.\ do LLM vs real} \\ \cmidrule(lr){2-3} \cmidrule(lr){4-5} \cmidrule(lr){6-7}
Regime & A & D & A & D & A & D \\
\midrule
zero-shot & 81,3 (78,0–85,3) & 66,7 & 34,7 (17,3–49,3) & 63,3 (62,7–64,0) & 32,4 (28,0–36,7) & 50,4 (50,0–50,7) \\
4-shot & 89,2 (84,2–94,5) & 89,1 (75,3–95,9) & 86,1 (68,0–94,0) & 81,0 (76,0–85,3) & 81,9 (67,8–91,8) & 81,4 (68,5–91,8) \\
10-shot & 94,4 (90,7–98,6) & 91,4 (80,0–97,1) & 93,4 (86,7–98,0) & 83,9 (77,3–86,0) & 89,4 (81,4–96,4) & 89,0 (81,4–97,1) \\
20-shot & 95,5 (92,3–97,7) & 92,6 (91,5–94,6) & 93,9 (89,3–98,0) & 84,9 (82,7–87,3) & 90,9 (84,6–95,4) & 90,4 (86,2–95,4) \\
40-shot & 95,2 (90,9–98,2) & 91,3 (88,2–92,7) & 93,4 (88,0–98,0) & 86,1 (84,7–88,0) & 91,4 (83,6–96,4) & 91,9 (87,3–96,4) \\
\midrule
W ajustado ao rótulo real (teto) & — & — & 96,2 (95,3–98,0) & 85,6 (84,7–86,7) & — & — \\
\bottomrule
\multicolumn{7}{p{0.95\linewidth}}{\footnotesize Perceptron, R$^2$. Mesmo protocolo nos dois problemas (150 pontos, exemplos do rótulo real, sem split). Fidelidade (vs LLM); acurácia (vs rótulo real) da métrica e do LLM como contexto; o teto (W ajustado ao rótulo real) só tem acurácia. Em zero-shot (e com nomes de classe neutros) a correspondência entre nome e classe é arbitrária: a acurácia (vs rótulo real) do LLM pode sair abaixo de 50\%; o valor invariante à permutação é max($a$, 1 $-$ $a$).} \\
\end{tabular}
\end{table}
